\documentclass[preview]{standalone}

\usepackage{amsmath}
\usepackage{amssymb}
\usepackage{stellar}
\usepackage{definitions}

\begin{document}

\id{psicologia-memoria-processi}
\genpage

\section{Memoria e processi mnestici}

\begin{snippet}{memoria-introduzione}
    La memoria permette di attribuire continuità alla propria esperienza e alla
    propria identità. Senza ricordi del passato sarebbe difficile dare senso a
    ciò che ci circonda; senza la capacità di acquisire nuovi ricordi, le
    esperienze recenti svanirebbero e diventerebbe difficile seguire una
    conversazione o una sequenza di eventi.
\end{snippet}

\begin{snippetdefinition}{memoria-definition}{Memoria}
    La \emph{memoria} è la capacità di codificare, immagazzinare e recuperare
    informazioni.
\end{snippetdefinition}

\begin{snippet}{funzioni-memoria}
    La memoria consente l'accesso consapevole agli eventi passati personali e
    collettivi, ma svolge anche funzioni implicite: rende familiari ambienti,
    oggetti e situazioni, permette di riconoscere regolarità e sostiene la
    continuità dell'esperienza quotidiana.
\end{snippet}

\includesnpt[width=48\%,src=/snippet/static-psychology/implicit-memory-out-of-context.jpg]{centered-img}

\begin{snippetdefinition}{uso-esplicito-memoria-definition}{Uso esplicito della memoria}
    L'\emph{uso esplicito della memoria} comporta uno sforzo consapevole per
    codificare o recuperare informazioni.
\end{snippetdefinition}

\begin{snippetdefinition}{uso-implicito-memoria-definition}{Uso implicito della memoria}
    L'\emph{uso implicito della memoria} si verifica quando la codifica o il
    recupero delle informazioni avvengono senza che il soggetto se ne renda
    conto.
\end{snippetdefinition}

\begin{snippetexample}{memoria-implicita-cucina-example}{Elemento fuori contesto}
    Se in una scena di cucina compare un coniglio, l'elemento appare
    immediatamente fuori luogo. Questa sensazione dipende dal confronto implicito
    tra la scena attuale e le conoscenze pregresse su ciò che normalmente si
    trova in una cucina, senza bisogno di un controllo consapevole di ogni
    oggetto presente.
\end{snippetexample}

\begin{snippetdefinition}{memoria-dichiarativa-definition}{Memoria dichiarativa}
    La \emph{memoria dichiarativa} riguarda il recupero di fatti o eventi che
    possono essere dichiarati verbalmente o rappresentati consapevolmente.
\end{snippetdefinition}

\begin{snippetdefinition}{memoria-procedurale-definition}{Memoria procedurale}
    La \emph{memoria procedurale} riguarda l'acquisizione, il mantenimento e
    l'uso delle abilità e delle procedure con cui si fanno le cose. È una forma
    implicita di memoria.
\end{snippetdefinition}

\begin{snippetdefinition}{codifica-definition}{Codifica}
    La \emph{codifica} è la fase iniziale dell'elaborazione mnestica che porta
    alla costruzione di una rappresentazione della memoria.
\end{snippetdefinition}

\begin{snippetdefinition}{immagazzinamento-definition}{Immagazzinamento}
    L'\emph{immagazzinamento} è il processo di conservazione nel tempo del
    materiale codificato.
\end{snippetdefinition}

\begin{snippetdefinition}{recupero-mnestico-definition}{Recupero mnestico}
    Il \emph{recupero mnestico} è il reperimento, in un momento successivo, di
    informazioni precedentemente codificate e immagazzinate.
\end{snippetdefinition}

\begin{snippet}{processi-mnestici}
    Codifica, immagazzinamento e recupero sono processi distinti ma
    interdipendenti: la codifica permette l'ingresso delle informazioni,
    l'immagazzinamento le conserva e il recupero le rende disponibili quando
    sono necessarie.
\end{snippet}

\begin{snippetdefinition}{rappresentazione-mentale-definition}{Rappresentazione mentale}
    Una \emph{rappresentazione mentale} è una rappresentazione interna che sta al
    posto di un oggetto o evento e trasmette informazioni congruenti con esso.
\end{snippetdefinition}

\begin{snippetexample}{rappresentare-oggetto-example}{Rappresentare un oggetto}
    Se si deve descrivere a un amico un regalo non presente nella stanza, si
    possono descriverne le proprietà, disegnarlo oppure mimarne l'uso. In tutti
    questi casi si costruiscono rappresentazioni interne che sostituiscono
    l'oggetto originario, conservandone alcune informazioni e tralasciandone
    altre.
\end{snippetexample}

\end{document}
